\section{Conclusion}

We introduced \textsc{iris}, an agent that learns purely in the imagination of a world model composed of a discrete autoencoder and an autoregressive Transformer. \textsc{iris} sets a new state of the art in the Atari 100k benchmark for methods without lookahead search. We showed that its world model acquires a deep understanding of game mechanics, resulting in pixel perfect predictions in some games. We also illustrated the generative capabilities of the world model, providing a rich gameplay experience when training in imagination. Ultimately, with minimal tuning compared to existing battle-hardened agents, \textsc{iris} opens a new path towards efficiently solving complex environments.

In the future, \textsc{iris} could be scaled up to computationally demanding and challenging tasks that would benefit from the speed of its world model. Besides, its policy currently learns from reconstructed frames, but it could probably leverage the internal representations of the world model. Another exciting avenue of research would be to combine learning in imagination with MCTS. Indeed, both approaches deliver impressive results, and their contributions to agent performance might be complementary.